\section{Which representation helps a class-requested journey?}
\label{sec:two-gafs}
One complete candidate provides two time series. Applying the spatial Hilbert mapping to its ordered positions yields the sequence for the inherited trajectory image. Applying a finite difference in physical time yields its speed sequence. For the latter we sum displacement magnitudes over complete one-second intervals and resample the resulting speed curve at 128 positions spanning the journey. A fixed 0--35 m/s amplitude scale preserves differences in physical speed. The speed sequence yields a summation angular field, a difference angular field and a transition field of discretized speeds. Figure~\ref{fig:speed-images} shows the first of those channels. A pixel at row $i$, column $j$ records a relation between speeds at two relative times; the axes do not represent geographic coordinates.

\fig{speed-images.pdf}{.96}{Two diffusion draws on the same 2.34 km bus route shown in Figure~\ref{fig:whole-opening}. Each image comes from the speed curve immediately to its left. The entire curve contributes to the angular image; the display colour scale is shared.}{speed-images}

The spatial CNN predicts participant class from the trajectory image alone. A separate CNN predicts six motion modes from the speed image alone. We use training speed tertiles (8.58 and 11.08 m/s) and divide each speed band by its training median four-second speed change. Given a class and a planned route, we sample a desired mode from the 64 nearest training trips by route length within that class. Class information therefore determines a distribution over desired motion modes; the speed-image reader itself has no class target. An alternative control omits the requested class when drawing the target. No evaluation reference is supplied to either selector.

We ran eight fresh same-map SUMO episodes and retained 3,021 arrival-confirmed requests. Two fixed diffusion weights each produced four candidates on the same prescribed route. We compared a uniform choice with selection by spatial class probability, selection by speed-mode probability and an equal-rank combination. Two further controls used a class-blind speed target and a direct calculation of the target mode from the candidate speeds. The spatial and speed readers achieved 25.1\% accuracy on eight classes and 82.9\% on six motion modes, respectively. These measure different tasks. The relevant head-to-head comparison is the quality of the delivered trajectories.

\begin{table}[htbp]\centering\small\setlength{\tabcolsep}{3pt}
\caption{Errors on fresh complete-trip traffic, averaged equally over participant classes. Lower is better. The first three columns are distributional Wasserstein discrepancies. The fourth is total variation in the motion-mode distribution.}
\begin{tabular}{lrrrr}\toprule
Selection & Steady run (s) & Trip time (s) & Change (m/s) & Mode TV\\\midrule
Uniform & 8.23 & 8.40 & .171 & .071\\
Spatial class GAF & 8.07 & 7.52 & .159 & .067\\
Speed-mode GAF & 8.59 & 8.72 & .149 & .074\\
Class-blind speed target & 8.28 & 8.22 & .166 & .059\\
Equal-rank combination & 8.53 & 8.25 & .153 & .067\\
Direct numeric mode & 9.41 & 9.79 & .151 & .084\\\bottomrule
\end{tabular}
\end{table}

Spatial feedback reduced travel-time discrepancy by 0.88 s relative to uniform selection (95\% paired episode interval 0.22 to 1.13 s) and by 1.21 s relative to class-directed speed feedback (0.44 to 1.51 s). Neither individual feedback strategy established a reduction in longest steady-run discrepancy relative to uniform. The speed selector reduced four-second speed-change discrepancy by 0.022 m/s versus uniform (0.006 to 0.039 m/s). The direct speed--spatial comparison on that secondary metric included zero. Conditioning the speed-mode target on class did not establish an improvement over the class-blind control on the two primary endpoints. This is evidence about terminal selection at four candidates per request in the specified demand model. Representation, training target and selection rule all differ between the two policies.

\subsection{Does speed guidance during sampling raise the spatial class score?}
Terminal selection can change which candidate is delivered without changing any candidate's motion. We therefore tested a different mechanism. During each reverse-diffusion step we predicted a clean travel time and progress profile, converted that estimate into a one-second along-road speed series, and differentiated the requested speed-mode log probability. The summation and difference angular channels supplied the gradient. The discretized transition channel contributed to the forward score only. A bounded noise correction with a smooth schedule acted throughout sampling. The target mode came from the training-only class and route-length prior. The spatial image classifier ran only after the final coordinates were decoded.

Eight new SUMO episodes supplied 192 requests, three per class per episode. Two generator seeds and four matched initial noises yielded 1,536 trajectories per arm. A control permuted the requested modes across participants, preserving their overall frequencies. Class-macro speed-mode probability increased from 0.364 to 0.380, a change of $+0.016$ (paired episode interval $+0.009$ to $+0.023$). The independently scored probability of the requested spatial class changed from 0.2370 to 0.2388, or $+0.19$ percentage points (interval $-0.24$ to $+0.65$ points). The permuted-mode control also reached 0.2388. The speed reader responded to its own guidance; transfer to spatial class recognition remained unresolved.

At the same four-candidate budget, terminal speed selection alone gave spatial probability 0.2469. Adding in-sampler guidance before the same selection gave 0.2353, a secondary change of $-1.15$ percentage points (interval $-2.35$ to $-0.05$ points). Without terminal selection, changes in travel-time, sustained cruise and four-second speed-change discrepancies all had paired intervals covering zero. The guided sampler took roughly 4.3 times as long as the unguided sampler in this run. These comparisons use one prespecified strength and two fixed generator weights. They do not establish a general improvement from speed guidance or from maximizing a participant classifier's score.

\subsection{Does the classification score hide shared motion information?}
Speed is derived from ordered positions and elapsed time. A speed image can therefore share information with the spatial trajectory image. The final spatial class probability reduces a 512-dimensional representation to one number. A weak change in that number alone cannot distinguish lack of shared information from lack of a useful change along the requested class direction. We examined the frozen readers' penultimate activations on complete SUMO journeys, keeping velocity features outside the spatial class reader.

We measured linear centered kernel alignment \cite{kornblith2019} on the first 32 training-fitted principal components of each representation. These components retained 81\% of spatial and 95\% of speed feature variance. On 3,021 observed G23 journeys the alignment was 0.573. Permuting the pairing of individual trips within the same class and traffic episode gave a median of 0.562. Thus much of the overall similarity survives the loss of individual trip correspondence. Across the six motorized classes alone, alignment was 0.011 (permutation median 0.004). After accounting linearly for class, route length, travel time, mean speed and endpoints, the all-class value was 0.0085 (permutation median 0.0083). The smaller G24 population likewise showed no clear excess over this conditional control. Adjustment also removes some genuine class-associated motion, so these residual comparisons cannot establish that information is absent from the original features. Alignment coefficients are not fractions of information shared.

A linear class probe trained on the frozen spatial embeddings, with regularization selected on separate monitor episodes, reached 25.1\% accuracy on the G23 journeys; the original head also reached 25.1\%. Under G24 guidance, the probe's requested-class probability increased by 0.66 percentage points (95\% paired interval $-0.06$ to $+1.57$), and the cosine margin to the requested class prototype changed by $-0.0006$ (interval $-0.0069$ to $+0.0058$). An empirical class-conditional kernel distance to reference spatial features was 0.01419 without guidance and 0.01433 with guidance. The difference interval covered zero. In the speed features the analogous distance increased from 0.02888 to 0.05099; almost all of the average increase came from pedestrians, while the six motorized classes had an inconclusive change. Distances are specific to the frozen features and the training-fitted kernel scale. Increasing a preferred mode's score does not guarantee reproduction of the entire observed class distribution.

Finally we fixed the same noisy diffusion state and road at 48 bus/car-compatible requests. At ten points on the unguided reverse path we compared the speed gradient with the change in the predicted clean state when switching the denoiser's class input between bus and private car. The schedule-weighted signed cosine averaged $0.0037$ (paired interval $-0.0022$ to $+0.0114$); its difference from a permuted mode target was $0.0011$ (interval $-0.0074$ to $+0.0092$). This probes a local direction in the model's normalized state and is specific to these two classes. It permits overlap at individual requests or noise levels. Taken together, the checks support broad shared structure while leaving the transfer of the current speed guidance to a more characteristic spatial class representation unestablished.

The paired-teacher study in Chapter~\ref{ch:paired} uses the same participant target for both encodings and transfers hidden features during diffusion training. It addresses a different hypothesis from the six-mode speed guidance and angular-value losses above.
